% Part: normal-modal-logic
% Chapter: filtrations
% Section: introduction

\documentclass[../../../include/open-logic-section]{subfiles}

\begin{document}

\olfileid{nml}{fil}{int}

\olsection{ভূমিকা}

কোনো যুক্তি সম্পর্কে একটি গুরুত্বপূর্ণ প্রশ্ন হলো সেটি সিদ্ধান্তযোগ্য কি না; অর্থাৎ, “এই !!{formula} কি সর্বত্র বৈধ?” প্রশ্নের উত্তর দেওয়ার কার্যকর পদ্ধতি আছে কি না। বচনযুক্তি সিদ্ধান্তযোগ্য: সত্যসারণি তৈরি করে কার্যকরভাবে পরীক্ষা করা যায় !!a{formula} টটোলজি কি না; প্রদত্ত !!{formula}-র সত্যসারণিটি সসীম। কিন্তু কোনো মোডাল সূত্র সব মডেলে সত্য কি না, তা সরাসরি পরীক্ষা করা যায় না, কারণ মডেলের সংখ্যা অসীম। একটি নির্দিষ্ট সূত্রের ক্ষেত্রে প্রাসঙ্গিক সব সসীম মডেল আমরা তালিকাভুক্ত করতে পারি: কেবল ওই !!{formula}-তে বাস্তবে উপস্থিত !!{propositional variable}s-এর জন্য জগৎসমষ্টির কোন কোন উপসমষ্টি নির্ধারিত হয়, সেটিই প্রাসঙ্গিক। অভিগম্যতা-সম্বন্ধ নির্দিষ্ট থাকলে $V(p)$-র সম্ভাব্য মানগুলি হলো~$W$-র সব উপসমষ্টি; $\card{W} = n$ হলে এমন উপসমষ্টি $2^n$টি। আমাদের !!{formula}~$!A$-তে $m$টি !!{propositional
  variable}s থাকলে $n$টি জগৎবিশিষ্ট ভিন্ন মডেল হয় $2^{nm}$টি। প্রতিটি মডেলে প্রতিটি জগতে~$!A$-র সত্যমান গণনা করেই পরীক্ষা করা যায় সব জগতে~$!A$ সত্য কি না। অবশ্য সম্ভাব্য সব অভিগম্যতা-সম্বন্ধও পরীক্ষা করতে হবে। কিন্তু $n$টি জগতে সম্বন্ধের সংখ্যাও সসীম: $W \times W$-র উপসমষ্টির সংখ্যা, অর্থাৎ $2^{n^2}$।

যদি আমাদের আগ্রহ \Log{K} যুক্তিতে না হয়ে কোনো বিশেষ শ্রেণির মডেল দিয়ে সংজ্ঞায়িত যুক্তিতে থাকে (যেমন প্রতিবিম্বী ও পরিযায়ী মডেলগুলির যুক্তি), তবে অভিগম্যতা-সম্বন্ধটি উপযুক্ত ধরনের কি না, সেটিও পরীক্ষা করতে হবে। যে ফ্রেমগুলি আমাদের আগ্রহের, সেগুলি মোডাল !!{formula}s দিয়ে সংজ্ঞায়িত হলে তা করা যায় (যেমন ফ্রেমে \Ax{T} ও \Ax{4} বৈধ কি না পরীক্ষা করে)। অতএব সব সসীম ফ্রেম একে একে নিয়ে, প্রত্যেকটি আমাদের কাঙ্ক্ষিত শ্রেণিতে পড়ে কি না পরীক্ষা করা, তারপর সেই ফ্রেমের সব সম্ভাব্য মডেল তালিকাভুক্ত করে প্রতিটিতে $!A$ সত্য কি না পরীক্ষা করাই ধারণাটি। সত্য না হলে থামব: কাঙ্ক্ষিত শ্রেণির মডেলগুলিতে $!A$ বৈধ নয়।

এই ধারণায় একটি সমস্যা আছে: কখন, আদৌ কখনো, খোঁজা থামানো যাবে তা আমরা জানি না। !!{formula}-টির কোনো সসীম প্রতিমডেল থাকলে আমাদের পদ্ধতি সেটি খুঁজে পাবে। কিন্তু সসীম প্রতিমডেল না থাকলে উত্তর মিলবে না। !!{formula}-টি বৈধ হতে পারে (তখন কোনো প্রতিমডেলই নেই), অথবা এর কেবল অসীম প্রতিমডেল থাকতে পারে, যা আমরা কখনো পরীক্ষা করব না। এই সমস্যা কাটানো যায় যদি দেখাতে পারি, প্রতিমডেলযুক্ত প্রতিটি !!{formula}-র একটি সসীম প্রতিমডেলও আছে। সে ক্ষেত্রে যুক্তিটির \emph{সসীম মডেল ধর্ম} আছে বলি।

কিন্তু কোনো যুক্তির সসীম মডেল ধর্ম আছে কীভাবে দেখাব? একটি উপায় হলো~$!A$-র একটি অসীম (প্রতি)মডেলকে সসীম মডেলে রূপান্তরের পদ্ধতি খোঁজা। তা সম্ভব হলে, যখনই কোনো মডেলে $!A$ সত্য নয়, প্রাপ্ত সসীম মডেলেও $!A$ সত্য হবে না। সব সসীম মডেলের তালিকায় ওই মডেলটি আসবে এবং অবৈধ প্রত্যেক সূত্রের অবৈধতা আমরা শেষ পর্যন্ত নির্ণয় করব। সূত্রটি বৈধ হলে এই পদ্ধতি থামবে না। যদি আরও দেখানো যায় যে প্রাপ্ত সসীম মডেলের আকারের একটি ঊর্ধ্বসীমা আছে এবং সেটি কেবল !!{formula}~$!A$-র উপর নির্ভর করে, তবে প্রতিমডেল খোঁজার জন্য ওই আকার পর্যন্ত সসীম মডেল পরীক্ষা করলেই চলবে। ততক্ষণে প্রতিমডেল না মিললে কোনো প্রতিমডেল নেই। তখন যেকোনো !!{formula}~$!A$-র জন্য “$!A$ কি বৈধ?” প্রশ্নটি পদ্ধতিটি বাস্তবেই সিদ্ধান্ত করতে পারবে।

অসীম গঠনকে সসীম করার একটি কার্যকর কৌশল হলো প্রাসঙ্গিক বিচারে একই রকম আচরণকারী গঠনের !!{element}s-কে “এক করে নেওয়া”। $\mModel{M}$-তে এমন আচরণের অসীম সংখ্যক জগৎ থাকলে, আমরা \emph{সব} জগৎকে সসীম সংখ্যক (প্রতিটি সম্ভবত অসীম) “শ্রেণি”-তে ভাগ করতে পারি। অন্য কথায়, জগৎসমষ্টিকে এমনভাবে বিভাগ করতে চাই যাতে বিভাগের সংখ্যা সসীম হয়, যদিও প্রতিটি বিভাগে অসীম সংখ্যক জগৎ থাকতে পারে। তারপর একটি নতুন মডেল~$\mModel{M^*}$ সংজ্ঞায়িত করব, যার জগৎ হবে ওই বিভাগগুলি। পুরোনো মডেলের সসীম সংখ্যক বিভাগ নতুন মডেলে সসীম সংখ্যক জগৎ দেয়; অর্থাৎ, একটি সসীম মডেল। $w$ যে বিভাগে আছে তাকে $[w]$ বলি। আমরা চাইব $\mSat{M}{!A}[w]$ তখনই এবং কেবল তখনই, যখন $\mSat{M^*}{!A}[{[w]}]$; কারণ পুরোনো মডেলটি $!A$-র প্রতিমডেল হলে নতুনটিকেও প্রতিমডেল হতে হবে। এর জন্য বিভাগ এবং~$\mModel{M^*}$-এর অভিগম্যতা-সম্বন্ধ যথাযথভাবে সংজ্ঞায়িত করতে হবে।

% BN-SRC-360: Interpret the box clause by universal accessibility and all-world truth.
% BN-SRC-361: Restore the valuation argument in V^*(p).
পদ্ধতিটি বোঝার জন্য প্রথমে এমন অবস্থা কল্পনা করি যেখানে অভিগম্যতা-সম্বন্ধ সর্বজনীন। তখন $\mSat{M}{\Box !B}[w]$ তখনই, যখন প্রত্যেক $v \in W$-তে $\mSat{M}{!B}[v]$; $\mModel{M^*}$-এর ক্ষেত্রেও একই কথা, শুধু $[w]$ ও $[v]$ নিয়ে। প্রথম প্রচেষ্টা হিসাবে ধরা যাক, দুটি জগৎ $u$ ও $v$ সমতুল্য (একই বিভাগে), যদি $\mModel{M}$-এর সব !!{propositional variable}s-এ তারা একমত হয়; অর্থাৎ, $\mSat{M}{p}[u]$ তখনই, যখন $\mSat{M}{p}[v]$। ধরি $V^*(p) = \Setabs{[w]}{\mSat{M}{p}[w]}$। আমাদের লক্ষ্য, $\mSat{M}{!A}[w]$ তখনই এবং কেবল তখনই, যখন $\mSat{M^*}{!A}[{[w]}]$। স্বাভাবিকভাবেই এটি আবেশে প্রমাণ করব। ভিত্তি হবে $!A \equiv p$। প্রথমে ধরি $\mSat{M}{p}[w]$। সংজ্ঞা থেকে $[w] \in V^*(p)$, তাই $\mSat{M^*}{p}[{[w]}]$। এবার ধরি $\mSat{M^*}{p}[{[w]}]$। তার মানে $[w] \in V^*(p)$; অর্থাৎ, $w$-এর সমতুল্য কোনো $v$-র জন্য $\mSat{M}{p}[v]$। কিন্তু “$w$ ও $v$ সমতুল্য” মানে “$w$ ও $v$-তে একই !!{propositional variable}s সত্য”; তাই $\mSat{M}{p}[w]$। এবার আবেশী ধাপ, যেমন $!A
\ident \lnot !B$। তখন $\mSat{M}{\lnot !B}[w]$ তখনই, যখন $\mSat/{M}{!B}[w]$; আবেশী অনুমান অনুসারে তখনই, যখন $\mSat/{M^*}{!B}[{[w]}]$; এবং তখনই, যখন $\mSat{M^*}{\lnot !B}[{[w]}]$। অন্য অ-মোডাল অপারেটরগুলির ক্ষেত্রেও একই যুক্তি চলে। $\Box$-এর ক্ষেত্রেও চলে: ধরি $\mSat{M^*}{\Box
  !B}[{[w]}]$। তার মানে প্রত্যেক $[u]$-তে $\mSat{M^*}{!B}[{[u]}]$। আবেশী অনুমানে প্রত্যেক $u$-তে $\mSat{M}{!B}[u]$। অতএব $\mSat{M}{\Box !B}[w]$।

সাধারণ ক্ষেত্রে $\mModel{M^*}$-এর অভিগম্যতা-সম্বন্ধও সংজ্ঞায়িত করতে হয়; তাই বিষয়টি আরও জটিল। ওপরের আবেশী প্রমাণটি কার্যকর রাখতে যে শর্তগুলি দরকার, অভিগম্যতা-সম্বন্ধ~$R^*$ সেগুলি মানলে মডেল~$\mModel{M^*}$-কে একটি \emph{ফিল্ট্রেশন} বলব। তখন যেকোনো ফিল্ট্রেশন~$\mModel{M^*}$-এ $[w]$-তে $!A$ সত্য তখনই, যখন $\mModel{M}$-এ $w$-তে $!A$ সত্য। তবে এমন ফিল্ট্রেশন যে \emph{আছে}, অর্থাৎ ওই শর্তগুলি মেনে~$R^*$ সংজ্ঞায়িত করা যায়, সেটিও দেখাতে হবে। এর জন্য $u$ ও $v$-কে কেবল সব !!{propositional variable}s-এ একমত হলে নয়, $!A$-র সব উপ-!!{formula}s-এ একমত হলেই সমতুল্য ধরতে হবে। $!A$-র উপ-!!{formula}s-এর সংখ্যা সসীম, তাই ফিল্ট্রেশনও সসীম হবে। ফিল্ট্রেশন সংজ্ঞায়নের উপায় একাধিক; তার অভিগম্যতা-সম্বন্ধে কাঙ্ক্ষিত ধর্ম (যেমন প্রতিবিম্বিতা, পরিযায়িতা ইত্যাদি) নিশ্চিত করতে~$R^*$-এর সংজ্ঞা বেছে নিতে হয়।

\end{document}
